\section{Movimiento rectilíneo uniforme (MRU)}

En el MRU, la velocidad es constante y la trayectoria es una línea recta.
\begin{align}
    v &= \frac{\Delta x}{\Delta t} = \text{constante} \\
    x &= x_0 + v \cdot t
\end{align}

\section{Movimiento rectilíneo uniformemente acelerado (MRUA)}

En el MRUA, la aceleración es constante.
\begin{align}
    v_f &= v_i + a \cdot t \\
    x &= x_i + v_i \cdot t + \frac{1}{2} a \cdot t^2 \\
    v_f^2 &= v_i^2 + 2a \cdot (x_f - x_i)
\end{align}

\subsection{Gráficas del movimiento}
\begin{itemize}
    \item \textbf{Posición vs. Tiempo:} Para MRU es una línea recta; para MRUA es una parábola.
    \item \textbf{Velocidad vs. Tiempo:} Para MRU es una línea horizontal; para MRUA es una línea recta con pendiente \(a\).
    \item \textbf{Aceleración vs. Tiempo:} Para MRU es cero; para MRUA es constante.
\end{itemize}

\section{Caída libre}

Es un caso particular de MRUA donde la aceleración es la gravedad (\(g = 9.8 \, \text{m/s}^2\)) y el movimiento es vertical.
\begin{equation}
    y = y_0 + v_0 t - \frac{1}{2} g t^2
\end{equation}

\resumebox{ejercicio}{Fuerza magnética del protón}{\footnotesize
    Un automóvil parte del reposo y acelera a \(2 \, \text{m/s}^2\) durante 10 segundos. ¿Qué distancia recorre? ¿Cuál es su velocidad final?
}